\begin{abstract}
Triton has emerged as the dominant intermediate language and compiler for deep learning primitives, providing the foundation for highly optimized libraries such as PyTorch 2.0's Inductor and FlashAttention. To achieve optimal performance, Triton provides a built-in autotuning mechanism (\texttt{@triton.autotune}) that performs an on-device search over kernel configurations (e.g., block sizes, number of warps, and pipeline stages). However, this autotuner exclusively optimizes for execution latency, entirely ignoring energy consumption---a critical metric given the escalating power demands of large language model (LLM) inference. In this paper, we present GreenTune, an energy-aware extension to Triton's autotuning framework. By introducing real-time power sampling via NVML directly into the benchmarking loop, GreenTune exposes the non-trivial Pareto frontier between latency and energy consumption. Across a suite of representative kernels including vector addition, Softmax, RMSNorm, General Matrix Multiplication (GEMM), and Fused Causal Attention on an NVIDIA Tesla T4 GPU, we demonstrate that the latency-optimal configuration is rarely energy-optimal. For instance, in memory-bound kernels such as RMSNorm, reducing the active warp count from 32 to 16 yields substantial energy savings with minimal latency degradation. Furthermore, we demonstrate that exhaustive grid searches for energy-optimal configurations can be effectively replaced by machine learning-guided multi-objective optimization using the Tree-structured Parzen Estimator (TPE), drastically reducing autotuning overhead. GreenTune demonstrates that energy can and should be treated as a first-class objective within modern tensor compilers.
\end{abstract}
\begin{IEEEkeywords}
GPU, Autotuning, Energy Efficiency, Triton, Deep Learning, Compilers
\end{IEEEkeywords}
